\chapter{Middens en parallellen}\label{ch:g8:midpoints}

Verbind de middens van twee zijden van een driehoek: het \omterm{def:g6:lines:objects}{lijnstuk} dat je krijgt
is altijd \omterm{def:g6:lines:perp}{parallel} met de derde zijde, en precies half zo lang. Deze
``middenparallelstelling'' en haar omgekeerde zijn de eerste voorsmaak van een
groot idee --- parallellen verdelen lengten in gelijke verhoudingen --- dat in
\cref{ch:g9:thales} uitgroeit tot de stelling van Thales.

\section{De middenparallelstelling}

\begin{theorem}[Middenparallelstelling]\label{thm:g8:midpoints:direct}
Zij in een driehoek $ABC$ het punt $I$ het midden van $[AB]$ en $J$ het midden
van $[AC]$. Dan is de \omterm{def:g6:lines:objects}{lijn} $(IJ)$ \omterm{def:g6:lines:perp}{parallel} met $(BC)$, en
\[
IJ = \frac{BC}{2} .
\]
\end{theorem}

\begin{proof}[Idee van het bewijs]
Zij $K$ het beeld van $I$ door de \omterm{def:g7:central:def}{puntspiegeling} om $J$ (een halve draai,
\cref{ch:g7:central}). In de \omterm{def:g4:geometry:polygon}{vierhoek} $AICK$ snijden de diagonalen $[IK]$ en
$[AC]$ elkaar in hun gemeenschappelijke midden $J$: het is dus een
\omterm{def:g7:central:parallelogram}{parallellogram}, zodat $KC$ \omterm{def:g6:lines:perp}{parallel} is met $AI$, dat wil zeggen met $(AB)$, en
$KC = AI = IB$. Dan heeft $IBCK$ twee overstaande zijden ($[IB]$ en $[KC]$) die
\omterm{def:g6:lines:perp}{parallel} en gelijk zijn: ook dat is een \omterm{def:g7:central:parallelogram}{parallellogram}
(\cref{thm:g7:central:recognize}), dus is $(IK) \parallel (BC)$ --- dat is
$(IJ) \parallel (BC)$ --- en $BC = IK = 2\,IJ$.
\end{proof}

\begin{omfigure}
\begin{tikzpicture}[scale=1.1]
  \coordinate (A) at (1.7,2.9);
  \coordinate (B) at (0,0);
  \coordinate (C) at (4.6,0);
  \coordinate (I) at ($(A)!0.5!(B)$);
  \coordinate (J) at ($(A)!0.5!(C)$);
  \draw[thick] (A) -- (B) -- (C) -- cycle;
  \draw[omThm, very thick] (I) -- (J);
  \fill (A) circle (1.5pt) node[above, font=\small] {$A$};
  \fill (B) circle (1.5pt) node[below left, font=\small] {$B$};
  \fill (C) circle (1.5pt) node[below right, font=\small] {$C$};
  \fill[omThm] (I) circle (1.5pt) node[left, font=\small] {$I$};
  \fill[omThm] (J) circle (1.5pt) node[right, font=\small] {$J$};
  \draw ($(A)!0.22!(B)$) ++(-0.09,-0.05) -- ++(0.18,0.1);
  \draw ($(A)!0.72!(B)$) ++(-0.09,-0.05) -- ++(0.18,0.1);
  \draw ($(A)!0.25!(C)$) ++(-0.06,-0.1) -- ++(0.12,0.2);
  \draw ($(A)!0.25!(C)$) ++(0.02,-0.1) -- ++(0.12,0.2);
  \draw ($(A)!0.75!(C)$) ++(-0.06,-0.1) -- ++(0.12,0.2);
  \draw ($(A)!0.75!(C)$) ++(0.02,-0.1) -- ++(0.12,0.2);
\end{tikzpicture}

{\small De middenparallel $[IJ]$ (rood) is \omterm{def:g6:lines:perp}{parallel} met de derde zijde $(BC)$ en
half zo lang.}
\end{omfigure}

\begin{example}\label{ex:g8:midpoints:direct}
In een driehoek $ABC$ met $BC = 9$ cm worden de middens $I$ van $[AB]$ en $J$ van
$[AC]$ verbonden. Zonder één meting: $(IJ) \parallel (BC)$ en
$IJ = \frac92 = 4.5$ cm. De kleine driehoek $AIJ$ is een op halve grootte
verkleinde kopie van $ABC$ --- ook zijn \omterm{def:g6:measure:perimeter}{omtrek} is de \omterm{def:g3:division:half}{helft} van de \omterm{def:g6:measure:perimeter}{omtrek} van
$ABC$.
\end{example}

\begin{theorem}[Omgekeerde stelling]\label{thm:g8:midpoints:converse}
In een driehoek $ABC$ snijdt de \omterm{def:g6:lines:objects}{lijn} door het midden $I$ van $[AB]$ die
\emph{\omterm{def:g6:lines:perp}{parallel}} is met $(BC)$, de zijde $[AC]$ in haar midden.
\end{theorem}
\admitted

\begin{example}\label{ex:g8:midpoints:converse}
Een pad steekt een driehoekig veld over, begint midden op de ene rand en loopt
\omterm{def:g6:lines:perp}{parallel} met de overstaande rand. De omgekeerde stelling waarborgt dat het pad
precies midden op de andere rand buitenkomt --- aan de overkant hoeft niets
gemeten te worden.
\end{example}

\begin{method}[Kiezen tussen stelling en omgekeerde]\label{met:g8:midpoints:choose}
\begin{enumerate}
  \item \emph{Twee middens bekend} $\Rightarrow$ gebruik de rechtstreekse
        stelling: besluit tot parallellisme en halve lengte.
  \item \emph{Eén midden en een \omterm{def:g6:lines:perp}{parallel} bekend} $\Rightarrow$ gebruik de
        omgekeerde: besluit dat het snijpunt een midden is.
  \item Schrijf op welke driehoek, welke middens en welke stelling je gebruikt
        --- telkens één zin.
\end{enumerate}
\end{method}

\section{Driehoeken op halve grootte}

\begin{proposition}[De middendriehoek]\label{prop:g8:midpoints:medial}
Het verbinden van de drie middens van de zijden van een driehoek verdeelt hem in
vier driehoeken met gelijke \omterm{def:g6:measure:area}{oppervlakte}, elk een op halve grootte verkleinde
kopie van de oorspronkelijke.
\end{proposition}

\begin{proof}
Elke zijde van de middendriehoek is \omterm{def:g6:lines:perp}{parallel} met een zijde van $ABC$ en half zo
lang (\cref{thm:g8:midpoints:direct}, drie keer toegepast). De vier kleine
driehoeken hebben zijden met diezelfde drie lengten, dus zijn het identieke
kopieën; samen betegelen ze $ABC$, dus heeft elk een \omterm{def:g3:division:half}{kwart} van zijn \omterm{def:g6:measure:area}{oppervlakte}.
\end{proof}

\begin{omfigure}
\begin{tikzpicture}[scale=1.05]
  \coordinate (A) at (1.9,3.0);
  \coordinate (B) at (0,0);
  \coordinate (C) at (5.0,0);
  \coordinate (I) at ($(A)!0.5!(B)$);
  \coordinate (J) at ($(A)!0.5!(C)$);
  \coordinate (K) at ($(B)!0.5!(C)$);
  \draw[thick] (A) -- (B) -- (C) -- cycle;
  \fill[omThm!12] (I) -- (J) -- (K) -- cycle;
  \draw[omThm, thick] (I) -- (J) -- (K) -- cycle;
  \fill (A) circle (1.4pt) node[above, font=\small] {$A$};
  \fill (B) circle (1.4pt) node[below left, font=\small] {$B$};
  \fill (C) circle (1.4pt) node[below right, font=\small] {$C$};
  \fill[omThm] (I) circle (1.4pt) node[left, font=\small] {$I$};
  \fill[omThm] (J) circle (1.4pt) node[right, font=\small] {$J$};
  \fill[omThm] (K) circle (1.4pt) node[below, font=\small] {$K$};
\end{tikzpicture}

{\small De middendriehoek $IJK$ verdeelt $ABC$ in vier gelijke driehoeken --- een
tekening om te onthouden.}
\end{omfigure}

\begin{remark}[Op weg naar Thales]
De middenparallelstelling is het geval ``één \omterm{def:g3:division:half}{helft}'' van een algemener feit: een
\omterm{def:g6:lines:objects}{lijn} \omterm{def:g6:lines:perp}{parallel} met één zijde van een driehoek verdeelt de twee andere zijden in
\emph{gelijke verhoudingen} --- een derde en een derde, twee vijfde en twee
vijfde\,\dots\ Die algemene uitspraak is de stelling van Thales
(\cref{ch:g9:thales}); dit jaar heb je alleen het geval van de middens nodig.
\end{remark}

\section{Oefeningen}

\begin{exercise}[$\star$]\label{exo:g8:midpoints:1}
In een driehoek $RST$ is $M$ het midden van $[RS]$ en $N$ het midden van $[RT]$,
met $ST = 7$ cm. Wat kun je zeggen over de \omterm{def:g6:lines:objects}{lijn} $(MN)$ en over de lengte $MN$?
Noem de gebruikte stelling.
\end{exercise}

\begin{exercise}[$\star$]\label{exo:g8:midpoints:2}
In een driehoek $ABC$ is $I$ het midden van $[AB]$, en de \omterm{def:g6:lines:objects}{lijn} door $I$ \omterm{def:g6:lines:perp}{parallel}
met $(BC)$ snijdt $[AC]$ in $J$; verder is $AC = 11$ cm. Bereken $AJ$ en noem de
gebruikte stelling.
\end{exercise}

\begin{exercise}[$\star$]\label{exo:g8:midpoints:3}
Teken een driehoek $ABC$ met $AB = 8$ cm, $AC = 6$ cm en $BC = 7$ cm, zet de
middens $I$ van $[AB]$ en $J$ van $[AC]$, en meet $IJ$ om de stelling na te gaan.
Bereken de \omterm{def:g6:measure:perimeter}{omtrek} van $AIJ$ zonder nog iets anders te meten.
\end{exercise}

\begin{exercise}[$\star$]\label{exo:g8:midpoints:4}
$IJK$ is de middendriehoek van $ABC$, waarvan de zijden $10$, $12$ en $16$ cm
meten. Geef de drie zijden van $IJK$ en zijn \omterm{def:g6:measure:perimeter}{omtrek}. Hoe verhouden de twee
omtrekken zich?
\end{exercise}

\begin{exercise}[$\star$]\label{exo:g8:midpoints:5}
De \omterm{def:g6:measure:area}{oppervlakte} van een driehoek is $36$ cm$^2$. Wat is de \omterm{def:g6:measure:area}{oppervlakte} van zijn
middendriehoek? En van elk van de vier kleine driehoeken?
\end{exercise}

\begin{exercise}[$\star\star$]\label{exo:g8:midpoints:6}
In een driehoek $ABC$ is $I$ het midden van $[AB]$, $J$ dat van $[AC]$ en $K$ dat
van $[BC]$. Toon aan dat $IJKB$ een \omterm{def:g7:central:parallelogram}{parallellogram} is. (Vergelijk $[IJ]$ en
$[BK]$: \omterm{def:g6:lines:perp}{parallel}? gelijk?)
\end{exercise}

\begin{exercise}[$\star\star$]\label{exo:g8:midpoints:7}
Een driehoekig zeil heeft een onderrand van $6.4$ m. Een versterkingsnaad
verbindt de middens van de twee andere randen. Hoeveel naad is er nodig? En als
de naad in de plaats daarvan de punten op een \omterm{def:g3:division:half}{kwart} van elke rand verbindt,
gerekend vanaf het bovenste \omterm{def:g5:solids:def}{hoekpunt}? (Vermoed het met een tekening; het bewijs
is dat van Thales, \cref{ch:g9:thales}.)
\end{exercise}

\begin{exercise}[$\star\star$]\label{exo:g8:midpoints:8}
In de driehoek $ABC$, \omterm{def:g6:shapes:triangles}{rechthoekig} in $A$, is $M$ het midden van de schuine zijde
$[BC]$, en de \omterm{def:g6:lines:perp}{parallel} met $(AB)$ door $M$ ontmoet $[AC]$ in $N$.
\begin{enumerate}
  \item Toon aan dat $N$ het midden van $[AC]$ is.
  \item Leg uit waarom $(MN)$ \omterm{def:g6:lines:perp}{loodrecht} op $(AC)$ staat.
\end{enumerate}
\end{exercise}

\begin{exercise}[$\star\star\star$]\label{exo:g8:midpoints:9}
Zij $ABCD$ een \emph{willekeurige} \omterm{def:g4:geometry:polygon}{vierhoek}, en $I$, $J$, $K$, $L$ de middens van
$[AB]$, $[BC]$, $[CD]$, $[DA]$ (het vermoeden van \cref{exo:g7:central:11}).
\begin{enumerate}
  \item Pas de middenparallelstelling toe in de driehoek $ABC$ op het \omterm{def:g6:lines:objects}{lijnstuk}
        $[IJ]$, en in de driehoek $ACD$ op het \omterm{def:g6:lines:objects}{lijnstuk} $[LK]$: vergelijk elk met
        de diagonaal $[AC]$.
  \item Besluit dat $IJKL$ altijd een \omterm{def:g7:central:parallelogram}{parallellogram} is.
\end{enumerate}
\end{exercise}

\section{Opgave: de drie zwaartelijnen en het zwaartepunt}

\begin{problem}[{Weekendopgave --- de zwaartelijnen van een driehoek gaan door
één punt, op twee derde van elk van hen}]
\label{pb:g8:midpoints:1}
Knip een driehoek uit stevig karton en hij balanceert, volkomen vlak, op een
punt van een potlood dat op één bijzonder punt staat: het
\emph{zwaartepunt}\index{zwaartepunt} van de driehoek. Deze opgave zoekt dat punt.
Een \emph{zwaartelijn}\index{zwaartelijn} van een driehoek is het \omterm{def:g6:lines:objects}{lijnstuk} dat een
\omterm{def:g5:solids:def}{hoekpunt} met het \emph{midden} van de overstaande zijde verbindt. Je zult
bewijzen --- met de middenparallelstelling, twee keer gebruikt op een onverwachte
manier --- dat de drie \omterm{pb:g8:midpoints:1}{zwaartelijnen} alle door één punt gaan, en je zult dat punt
precies aanwijzen.

\medskip
\noindent\textbf{Deel I --- Opwarming.}
\begin{enumerate}
  \item Teken een grote driehoek $ABC$ (geen bijzondere vorm), construeer de
        middens $I$ van $[AB]$, $J$ van $[AC]$ en $K$ van $[BC]$, en teken de
        drie \omterm{pb:g8:midpoints:1}{zwaartelijnen} $[AK]$, $[BJ]$, $[CI]$. Wat neem je waar?
  \item Bewijs dat een \omterm{pb:g8:midpoints:1}{zwaartelijn} de driehoek in twee driehoeken met gelijke
        \omterm{def:g6:measure:area}{oppervlakte} verdeelt. (Vergelijk basissen en hoogten, en gebruik de
        formule voor de \omterm{def:g6:measure:area}{oppervlakte}: basis $\times$ hoogte $\div\ 2$.) Wat zijn
        voor een driehoek met \omterm{def:g6:measure:area}{oppervlakte} $36$ cm$^2$ de \omterm{def:g6:measure:area}{oppervlakten} van de twee
        stukken?
  \item Wat zegt \cref{thm:g8:midpoints:direct} in de driehoek $ABC$ over het
        \omterm{def:g6:lines:objects}{lijnstuk} $[IJ]$?
\end{enumerate}

\medskip
\noindent\textbf{Deel II --- Twee \omterm{pb:g8:midpoints:1}{zwaartelijnen} snijden elkaar op twee derde.}
De \omterm{pb:g8:midpoints:1}{zwaartelijnen} $[BJ]$ en $[CI]$ snijden elkaar in een punt; noem het $G$. Zij
$M$ het midden van $[GB]$ en $N$ het midden van $[GC]$.
\begin{enumerate}[resume]
  \item Pas de middenparallelstelling toe \emph{in de driehoek $GBC$}: wat zegt
        ze over het \omterm{def:g6:lines:objects}{lijnstuk} $[MN]$?
  \item Leid af dat $(IJ) \parallel (MN)$ en $IJ = MN$, en besluit met
        \cref{thm:g7:central:recognize} dat $IJNM$ een \omterm{def:g7:central:parallelogram}{parallellogram} is.
  \item Leg uit waarom $I$ en $N$ beide op de \omterm{pb:g8:midpoints:1}{zwaartelijn} $(CI)$ liggen, en $J$
        en $M$ beide op de \omterm{pb:g8:midpoints:1}{zwaartelijn} $(BJ)$ --- zodat de diagonalen van het
        \omterm{def:g7:central:parallelogram}{parallellogram} $IJNM$ de lijnstukken $[IN]$ en $[JM]$ zijn, en ze elkaar
        precies in $G$ snijden. Wat zegt de \emph{definitie} van een
        \omterm{def:g7:central:parallelogram}{parallellogram} dan over $G$?
  \item Leid af dat $BM = MG = GJ$, en besluit:
        \[
        BG = 2 \times GJ
        \]
        --- het punt $G$ ligt op de \omterm{pb:g8:midpoints:1}{zwaartelijn} $[BJ]$ op twee derde van de weg
        vanaf het \omterm{def:g5:solids:def}{hoekpunt} $B$. Formuleer en verantwoord het overeenkomstige
        resultaat op de \omterm{pb:g8:midpoints:1}{zwaartelijn} $[CI]$.
  \item Vergeet nu $[CI]$ en voer precies hetzelfde betoog uit voor het paar
        \omterm{pb:g8:midpoints:1}{zwaartelijnen} $[BJ]$ en $[AK]$: ook hun snijpunt ligt op twee derde van
        $[BJ]$ vanaf $B$. Leg uit waarom dat dwingt dat het \emph{hetzelfde} punt
        $G$ is --- en besluit dat alle drie de \omterm{pb:g8:midpoints:1}{zwaartelijnen} door $G$ gaan.
  \item Hoe ver ligt $G$ van $B$, en van $J$, in een driehoek waarvan de
        \omterm{pb:g8:midpoints:1}{zwaartelijn} $[BJ]$ $9$ cm meet?
  \item Verantwoord in één zin dat ook $AG = 2 \times GK$.
\end{enumerate}

\medskip
\noindent\textbf{Deel III --- Zes gelijke stukken, en coördinaten.}
De drie \omterm{pb:g8:midpoints:1}{zwaartelijnen} verdelen de driehoek in zes kleine driehoeken rond $G$.
\begin{enumerate}[resume]
  \item In de driehoek $GBC$ is het \omterm{def:g6:lines:objects}{lijnstuk} $[GK]$ een \omterm{pb:g8:midpoints:1}{zwaartelijn}. Leid af dat
        de twee kleine driehoeken $GBK$ en $GKC$ gelijke \omterm{def:g6:measure:area}{oppervlakte} hebben.
  \item Toon met vraag 2 in de driehoeken $ABK$ en $AKC$ (beide verdeeld door de
        \omterm{pb:g8:midpoints:1}{zwaartelijn} $[AK]$ van $ABC$) aan dat de driehoeken $ABG$ en $ACG$
        gelijke \omterm{def:g6:measure:area}{oppervlakte} hebben. Herhaal met een andere \omterm{pb:g8:midpoints:1}{zwaartelijn} en besluit
        dat de drie driehoeken $ABG$, $BCG$, $CAG$ alle een \omterm{def:g6:measure:area}{oppervlakte} hebben
        die een derde is van die van $ABC$.
  \item Elk van deze drie driehoeken wordt door een stuk \omterm{pb:g8:midpoints:1}{zwaartelijn} in twee
        gelijke \omterm{def:g3:division:half}{helften} verdeeld ($[GI]$ is bijvoorbeeld een \omterm{pb:g8:midpoints:1}{zwaartelijn} van de
        driehoek $ABG$ --- vanuit welk \omterm{def:g5:solids:def}{hoekpunt} gezien?). Besluit: de zes kleine
        stukken rond $G$ hebben alle dezelfde \omterm{def:g6:measure:area}{oppervlakte}, een zesde van die van
        de driehoek.
  \item Zet in een assenstelsel de punten $A(0, 0)$, $B(6, 0)$ en $C(0, 6)$, met
        $K(3, 3)$ het midden van $[BC]$. Bereken met de stelling over twee derde
        uit deel II de coördinaten van $G$ op de \omterm{pb:g8:midpoints:1}{zwaartelijn} $[AK]$; ga daarna na
        dat hetzelfde punt op twee derde van de \omterm{pb:g8:midpoints:1}{zwaartelijn} $[BJ]$ ligt, met
        $J(0, 3)$. Vergelijk de coördinaten van $G$ met de gemiddelden
        \[
        \frac{0 + 6 + 0}{3},
        \qquad
        \frac{0 + 0 + 6}{3} .
        \]
  \item Slot over het evenwicht: leg met vragen 2 en 8 uit waarom de kartonnen
        driehoek balanceert op een mesrand die langs een willekeurige \omterm{pb:g8:midpoints:1}{zwaartelijn}
        wordt gelegd --- aan elke kant evenveel karton --- en waarom de
        potloodpunt dus in $G$ moet staan, het punt dat de fysici het \omterm{pb:g8:midpoints:1}{zwaartepunt}
        noemen. (Een volkomen sluitend bewijs van dat evenwicht vraagt de
        integraalrekening uit de universitaire \omterm{def:g6:measure:volume}{volumes}; de gelijke \omterm{def:g6:measure:area}{oppervlakten}
        maken het vandaag geloofwaardig.)
\end{enumerate}
\end{problem}
